\documentclass[12pt,a4paper]{article}

\usepackage[margin=1in]{geometry}
\usepackage{graphicx}
\usepackage{booktabs}
\usepackage{float}
\usepackage{hyperref}
\usepackage{enumitem}
\usepackage{setspace}

\hypersetup{
    colorlinks=false,
    pdfborder={0 0 0},
    hidelinks
}

\setlength{\parindent}{0pt}
\setlength{\parskip}{6pt}
\onehalfspacing

\title{\textbf{CSE638 GRS 2026}\\[0.5em]
\textbf{Performance Comparison of Kubernetes CNI Plugins}\\
\large Flannel vs Calico vs Cilium\\[0.6em]
\normalsize Group Number: 5\\
Ayush Jain - MT25066\\
Singh Tharun - MT25087\\
Ruchir Jain - MT25080}
\author{}
\date{}

\begin{document}
\maketitle
\tableofcontents
\newpage

\section{Problem Statement}
In Kubernetes, pod networking depends on the CNI plugin. Each plugin works in its own way. Because of that, the same app can show different latency and different CPU/memory usage.

In this project, we compare three popular CNI plugins:
\begin{itemize}[leftmargin=1.5em]
    \item Flannel
    \item Calico
    \item Cilium
\end{itemize}

So we ran the same workload on all three and compared the output.

\section{Why This Problem Matters}
For production systems, networking affects:
\begin{itemize}[leftmargin=1.5em]
    \item User response time
    \item Cost of running the cluster
    \item Stability during high traffic
\end{itemize}

So the CNI choice is not just a setup step, it changes real performance. Our goal was to make this choice easier using actual measurements.

\section{Objectives}
Our goals were:
\begin{enumerate}[leftmargin=1.8em]
    \item Build the same benchmark setup for all three CNIs.
    \item Measure latency carefully, especially tail latency (p99).
    \item Measure CPU and memory overhead.
    \item Explain results in simple language.
\end{enumerate}

\section{Background (Short)}
\subsection{Flannel}
Flannel is simple and lightweight. In many cases, it is easy to manage and uses fewer resources.

\subsection{Calico}
Calico is strong for policy and routing. It usually gives a good balance of features and performance.

\subsection{Cilium}
Cilium uses eBPF. It gives advanced networking features and better visibility, but it may use more resources in some setups.

\section{Methodology}
\subsection{What We Kept Constant}
To keep the comparison fair, we kept these things same in every run:
\begin{itemize}[leftmargin=1.5em]
    \item Machine and environment
    \item Workload type
    \item Deployment process
    \item Measurement method
\end{itemize}

Only one thing changed: the CNI plugin.

\subsection{Metrics Collected}
\begin{itemize}[leftmargin=1.5em]
    \item Latency: p50, p90, p99, average
    \item Throughput and error rate
    \item Resource usage: CPU (mcores) and Memory (MiB)
\end{itemize}

\section{Installation and Run Steps}
From repository root:

\begin{verbatim}
chmod +x G_5_scripts/*.sh
./G_5_scripts/G_5_setup_prereqs.sh
\end{verbatim}

For each CNI (flannel/calico/cilium):

\begin{verbatim}
./G_5_scripts/G_5_create_cluster.sh <cni>
./G_5_scripts/G_5_deploy_workload.sh <cni>
./G_5_scripts/G_5_run_benchmark.sh <cni>
./G_5_scripts/G_5_collect_metrics.sh <cni>
\end{verbatim}

Plot generation:

\begin{verbatim}
./G_5_artifacts/G_5_venv/bin/python G_5_plots/G_5_plot_latency_hardcoded.py
./G_5_artifacts/G_5_venv/bin/python G_5_plots/G_5_plot_latency_variability_hardcoded.py
./G_5_artifacts/G_5_venv/bin/python G_5_plots/G_5_plot_overhead_hardcoded.py
\end{verbatim}

\section{Results and Observations}
\subsection{p99 Latency Comparison}
\begin{figure}[H]
    \centering
    \includegraphics[width=0.85\textwidth]{../G_5_plots/G_5_p99_latency.png}
    \caption{p99 latency comparison for Flannel, Calico, and Cilium.}
\end{figure}

\textbf{What we saw:}
All three CNIs were close in latency. The differences were small in our test range.

\subsection{Latency Variability Across Runs}
\begin{figure}[H]
    \centering
    \includegraphics[width=0.85\textwidth]{../G_5_plots/G_5_latency_variability.png}
    \caption{Mean p99 latency with min-max spread across repeated runs.}
\end{figure}

\textbf{What we saw:}
Results were mostly stable across repeated runs. This gave us confidence that the trend is real.

\subsection{CPU and Memory Overhead}
\begin{figure}[H]
    \centering
    \includegraphics[width=0.85\textwidth]{../G_5_plots/G_5_resource_overhead.png}
    \caption{Aggregate CPU and memory overhead.}
\end{figure}

\textbf{What we saw:}
Flannel had the lowest overhead in our setup. Calico stayed in the middle. Cilium used more CPU and memory.

\section{Why We Focused on p99}
Average latency can look good even when some requests are very slow. In real apps, users notice these slow requests.

p99 helps us capture that slow tail clearly.

So we kept p50, p90, and average for context, but used p99 as the main reliability indicator.

\section{Sample Data from CSV Files}
\subsection{Top 5 Rows from Latency CSVs (QPS 100)}

\begin{table}[H]
\centering
\small
\caption{Flannel latency sample rows.}
\begin{tabular}{cccccc}
\toprule
Run & p50 & p90 & p99 & avg & throughput \\
\midrule
1 & 0.00054596 & 0.00092766 & 0.00159493 & 0.00059488 & 99.9946 \\
2 & 0.00056824 & 0.00094095 & 0.00174056 & 0.00062646 & 99.9945 \\
3 & 0.00052654 & 0.00092067 & 0.00128243 & 0.00057646 & 99.9954 \\
4 & 0.00055438 & 0.00093574 & 0.00158489 & 0.00059925 & 99.9948 \\
5 & 0.00056222 & 0.00093156 & 0.00149200 & 0.00059302 & 99.9937 \\
\bottomrule
\end{tabular}
\end{table}

\begin{table}[H]
\centering
\small
\caption{Calico latency sample rows.}
\begin{tabular}{cccccc}
\toprule
Run & p50 & p90 & p99 & avg & throughput \\
\midrule
1 & 0.00055235 & 0.00094075 & 0.00174991 & 0.00060487 & 99.9935 \\
2 & 0.00055305 & 0.00094226 & 0.00136714 & 0.00059748 & 99.9923 \\
3 & 0.00055499 & 0.00094422 & 0.00177175 & 0.00061223 & 99.9945 \\
4 & 0.00055118 & 0.00092480 & 0.00120133 & 0.00055425 & 99.9953 \\
5 & 0.00055501 & 0.00092518 & 0.00130327 & 0.00056306 & 99.9933 \\
\bottomrule
\end{tabular}
\end{table}

\begin{table}[H]
\centering
\small
\caption{Cilium latency sample rows.}
\begin{tabular}{cccccc}
\toprule
Run & p50 & p90 & p99 & avg & throughput \\
\midrule
1 & 0.00055781 & 0.00094486 & 0.00179226 & 0.00066001 & 99.9953 \\
2 & 0.00057691 & 0.00096374 & 0.00186320 & 0.00068801 & 99.9927 \\
3 & 0.00056278 & 0.00095677 & 0.00145410 & 0.00067041 & 99.9944 \\
4 & 0.00057342 & 0.00094929 & 0.00143620 & 0.00064682 & 99.9956 \\
5 & 0.00055447 & 0.00095056 & 0.00156517 & 0.00064484 & 99.9932 \\
\bottomrule
\end{tabular}
\end{table}

\subsection{Top 5 Rows from Resource CSVs}
\begin{table}[H]
\centering
\small
\caption{Flannel resource sample rows.}
\begin{tabular}{llll}
\toprule
Namespace & Pod & CPU (mcores) & Memory (MiB) \\
\midrule
g5-bench & g5-client-d57c4d46-6q6rl & 1 & 5 \\
g5-bench & g5-echo-649fbccd57-7t4wj & 0 & 10 \\
g5-bench & g5-echo-649fbccd57-d7dmb & 0 & 10 \\
g5-bench & g5-echo-649fbccd57-dj2fn & 0 & 10 \\
kube-flannel & kube-flannel-ds-nfcnd & 9 & 14 \\
\bottomrule
\end{tabular}
\end{table}

\begin{table}[H]
\centering
\small
\caption{Calico resource sample rows.}
\begin{tabular}{llll}
\toprule
Namespace & Pod & CPU (mcores) & Memory (MiB) \\
\midrule
g5-bench & g5-client-d57c4d46-n6j2q & 0 & 4 \\
g5-bench & g5-echo-649fbccd57-ck6rc & 0 & 11 \\
g5-bench & g5-echo-649fbccd57-pztgs & 0 & 10 \\
g5-bench & g5-echo-649fbccd57-xxsnk & 0 & 10 \\
kube-system & calico-kube-controllers-5b4bd99c9c-ln7j6 & 4 & 13 \\
\bottomrule
\end{tabular}
\end{table}

\begin{table}[H]
\centering
\small
\caption{Cilium resource sample rows.}
\begin{tabular}{llll}
\toprule
Namespace & Pod & CPU (mcores) & Memory (MiB) \\
\midrule
g5-bench & g5-client-d57c4d46-9kwz7 & 1 & 4 \\
g5-bench & g5-echo-649fbccd57-42f7m & 0 & 8 \\
g5-bench & g5-echo-649fbccd57-8jvjg & 0 & 9 \\
g5-bench & g5-echo-649fbccd57-qrlc5 & 0 & 8 \\
kube-system & cilium-envoy-7j8pv & 11 & 23 \\
\bottomrule
\end{tabular}
\end{table}

\section{Conclusion}
This comparison gave us a clear and practical picture:
\begin{itemize}[leftmargin=1.5em]
    \item All three CNIs gave solid latency performance in our tested load range.
    \item The biggest difference was resource overhead.
    \item Flannel was most lightweight in this setup.
    \item Calico gave a balanced profile.
    \item Cilium offered strong features with higher resource cost.
\end{itemize}

In short, the best CNI depends on what we care about most: low overhead, balanced performance with policy, or advanced eBPF-based features.

\section{References}
\begin{enumerate}[leftmargin=1.8em]
    \item Kubernetes Documentation: \url{https://kubernetes.io/docs/}
    \item CNI Specification: \url{https://github.com/containernetworking/cni}
    \item Flannel Documentation: \url{https://github.com/flannel-io/flannel}
    \item Calico Documentation: \url{https://docs.tigera.io/calico/latest}
    \item Cilium Documentation: \url{https://docs.cilium.io/}
    \item Fortio Documentation: \url{https://github.com/fortio/fortio}
\end{enumerate}

\end{document}
